\documentclass[UTF8,fontset=fandol,zihao=5]{ctexart}

\usepackage[a4paper,margin=2cm]{geometry}
\usepackage{amsmath,amssymb,mathtools}
\usepackage{upgreek}
\usepackage[nolimits]{cmupint}
\usepackage{booktabs,tabularx,array}
\usepackage{xcolor}
\usepackage[most]{tcolorbox}
\usepackage{enumitem}
\usepackage{indentfirst}
\usepackage{hyperref}
\usepackage{tikz}
\usepackage{pdflscape}

\newfontfamily\russianfont{Times New Roman}

\usetikzlibrary{
  positioning,
  calc,
  decorations.pathreplacing,
  arrows.meta
}

% relationbox 中知识关系图的统一样式
\tikzset{
  relnode/.style={
    font=\small,
    align=center
  },
  relline/.style={
    thick
  },
  relarrow/.style={
    -{Stealth},
    thick
  }
}

\DeclareTextFontCommand{\textbf}{\heiti\color{main}}

\definecolor{main}{HTML}{B33F46}
\definecolor{mainlight}{HTML}{FFF0F0}
\definecolor{method}{HTML}{B3797C}
\definecolor{methodlight}{HTML}{FFF0F0}

\hypersetup{colorlinks=true,linkcolor=main,urlcolor=main}
\setlength{\parindent}{2em}
\linespread{1.5}
\setlength{\abovedisplayskip}{5pt}
\setlength{\belowdisplayskip}{5pt}
\setlength{\abovedisplayshortskip}{3pt}
\setlength{\belowdisplayshortskip}{3pt}
\setlist[itemize]{leftmargin=2em,itemsep=0.1em,topsep=0.2em}
\setlist[enumerate]{leftmargin=2.3em,itemsep=0.1em,topsep=0.2em}

\ctexset{
  section={format=\centering\heiti\zihao{3}\color{main},
            name={第,章},number=\chinese{section},aftername=\quad,
            beforeskip=0.9em,afterskip=0.5em},
  subsection={format=\heiti\zihao{5}\color{main},
              name={,、},number=\chinese{subsection},aftername=,
              beforeskip=0.6em,afterskip=0.22em},
  subsubsection={format=\heiti\zihao{5}\color{main},
                  number={（\chinese{subsubsection}）},aftername=,
                  beforeskip=0.55em,afterskip=0.15em}
}

\newtcolorbox{relationbox}{
  enhanced,breakable,colback=mainlight,colframe=main,
  boxrule=0.65pt,arc=1.4mm,left=2mm,right=2mm,top=1.2mm,bottom=1.2mm
}
\newtcolorbox{methodbox}[1][]{
  enhanced,breakable,colback=methodlight,colframe=method,
  boxrule=0.65pt,arc=1.4mm,left=2mm,right=2mm,top=1.2mm,bottom=1.2mm,
  title=#1,fonttitle=\kaishu,fontupper=\kaishu,
  before upper={\setlength{\parindent}{2em}}
}

\newcommand{\ii}{\mathrm{i}}
\newcommand{\ee}{\mathrm{e}}
\newcommand{\CC}{\mathbb{C}}
\newcommand{\RR}{\mathbb{R}}
\newcommand{\dd}{\,\mathrm{d}}
\newcommand{\FF}{\mathcal{F}}

\begin{document}

\begin{center}
  {\heiti\zihao{1}\color{main} 概率论与数理统计笔记}
\end{center}

\begin{relationbox}
\[
  \boxed{\text{概率论}}
  \Longleftrightarrow
  \boxed{\text{数理统计}}
\]
\end{relationbox}

\section{概率论}

\subsection{\russianfont Колмогоров公理化概率论}

\subsubsection{样本空间$\varOmega$}

\subsubsection{事件域$\FF$}

\subsubsection{概率$P$}

\subsection{常见概型}

\subsubsection{古典概型}

\subsubsection{几何概型}

\subsection{随机变量$X$}

\subsection{常见分布}

\subsubsection{离散型分布}

\begin{center}
  \begin{tabularx}{0.92\textwidth}{>{\bfseries}p{2.3cm}c c c c}
    \toprule
    名称 & 分布律 & 累积概率函数 & 期望$E(X)$ & 标准差$D(X)$ \\
    \midrule
    两点分布（0-1分布） & $\displaystyle P(X=1)=p,\ P(X=0)=1-p$ & $F(X)=P(X\leq x)=1-p+px$ & $E(X)=p$ & $\sqrt{p(1-p)}$ \\
    二项分布 & $\displaystyle P(X=k)=\binom{n}{k}p^k(1-p)^{n-k}$ & $F(X)=P(X\leq x)$ & $E(X)=np$ & $\sqrt{np(1-p)}$ \\
    Poisson分布 & $\displaystyle P(X=k)=\frac{\lambda^k e^{-\lambda}}{k!}$ & $F(X)=P(X\leq x)$ & $E(X)=\lambda$ & $\sqrt{\lambda}$ \\
    几何分布 & $\displaystyle P(X=k)=(1-p)^{k-1}p$ & $F(X)=P(X\leq x)$ & $E(X)=\frac{1}{p}$ & $\sqrt{\frac{1-p}{p^2}}$ \\
    Pascal分布 & $\displaystyle P(X=k)=\binom{k-1}{r-1}p^r(1-p)^{k-r}$ & $F(X)=P(X\leq x)$ & $E(X)=\frac{r}{p}$ & $\sqrt{\frac{r(1-p)}{p^2}}$ \\
    超几何分布 & $\displaystyle P(X=k)=\frac{\binom{M}{k}\binom{N-M}{n-k}}{\binom{N}{n}}$ & $F(X)=P(X\leq x)$ & $E(X)=\frac{nM}{N}$ & $\sqrt{\frac{nM(N-M)(N-n)}{N^2(N-1)}}$ \\
    \bottomrule
  \end{tabularx}
\end{center}

\subsubsection{连续型分布}

\begin{center}
  \begin{tabularx}{0.92\textwidth}{>{\bfseries}p{2.3cm}c c c c}
    \toprule
    名称 & 概率密度函数$f(x)$ & 累积概率函数$F(x)$ & 期望$E(X)$ & 标准差$D(X)$ \\
    \midrule
    均匀分布 & $\displaystyle f(x)=\frac{1}{b-a},\ a<x<b$ & $F(X)=P(X\leq x)=\frac{x-a}{b-a}$ & $E(X)=\frac{a+b}{2}$ & $\sqrt{\frac{(b-a)^2}{12}}$ \\
    指数分布 & $\displaystyle f(x)=\lambda e^{-\lambda x},\ x>0$ & $F(X)=P(X\leq x)=1-e^{-\lambda x}$ & $E(X)=\frac{1}{\lambda}$ & $\sqrt{\frac{1}{\lambda^2}}$ \\
    正态分布（Gauss分布） & $\displaystyle \varphi(x)=\frac{1}{\sigma \sqrt{2 \uppi}} e^{-\frac{(x-\mu)^2}{2 \sigma^2}}$ & $\varPhi(X)=P(X\leq x)$ & $E(X)=\mu$ & $D(X)=\sigma$ \\
    \bottomrule
  \end{tabularx}
\end{center}

\section{概率统计}

\subsection{频率}

\end{document}
